\hypertarget{news-rag-pipeline-on-aws}{%
\section{News RAG Pipeline on AWS}}

\hypertarget{a-serverless-retrieval-augmented-generation-pipeline-for-intelligent-news-qa}{%
\subsection{A Serverless Retrieval-Augmented Generation Pipeline for
Intelligent News
Q\&A}}

\hypertarget{executive-summary}{%
\subsubsection{1. Executive Summary}}

The News RAG Pipeline is designed to build an intelligent news Q\&A
system that automatically crawls news articles from Vietnamese news
sites (VnExpress, Thanh Nien, VietnamNet), processes them through a Data
Warehouse (Star Schema), generates vector embeddings using Amazon
Bedrock Titan Embed v2, and enables natural language querying through
Retrieval-Augmented Generation (RAG) --- entirely serverless on AWS. The
system serves as a practical demonstration of integrating AWS serverless
services with LLM APIs (Groq, Gemini) to create an AI-powered
information retrieval platform.

\hypertarget{problem-statement}{%
\subsubsection{2. Problem Statement}}

\hypertarget{whats-the-problem}{%
\paragraph{What's the Problem?}}

Keeping up with news from multiple sources requires manually reading and
searching through countless articles. There is no centralized system
that allows users to ask questions about recent news and get
AI-generated answers with proper citations. Existing solutions like
ChatGPT lack up-to-date news context without manual prompting.

\hypertarget{the-solution}{%
\paragraph{The Solution}}

The News RAG Pipeline automates the entire workflow: (1) Scrapy crawlers
on ECS Fargate collect articles from Vietnamese news sitemaps, (2)
articles are ingested via SQS into Aurora PostgreSQL, (3) Lambda ETL
processes raw HTML into a Star Schema and generates vector embeddings
via Amazon Bedrock Titan Embed v2, (4) a Lambda RAG API accepts natural
language queries, performs vector similarity search on pgvector (HNSW
index), and generates answers using Groq (Qwen3-8B, Llama 3.1) or Gemini
2.0 Flash LLMs --- all orchestrated by EventBridge Scheduler.

\hypertarget{benefits-and-return-on-investment}{%
\paragraph{Benefits and Return on
Investment}}

The solution provides a foundational platform for learning modern AWS
serverless architecture, RAG systems, and MLOps practices. Key benefits
include: automated news aggregation eliminates manual searching,
AI-powered Q\&A with source citations saves research time, hands-on
experience with 10+ AWS services, and cost-effective serverless
infrastructure. Monthly costs are approximately \$21-26 USD. The
break-even period is achieved through significant time savings from
automated information retrieval versus manual news monitoring.

\hypertarget{solution-architecture}{%
\subsubsection{3. Solution Architecture}}

The platform employs a serverless AWS architecture based on two pipeline
stages: (1) Data Pipeline --- EventBridge Scheduler triggers ECS Fargate
crawler daily at 01:00 UTC, pushing articles to SQS; Lambda Consumer
processes messages and inserts raw articles into Aurora PostgreSQL with
SHA256 deduplication. (2) ETL + RAG Pipeline --- second EventBridge
trigger at 02:00 UTC runs Lambda ETL that cleans HTML, chunks text at
500 tokens, generates 1024-dimensional embeddings via Bedrock Titan
Embed v2, and stores vectors in Aurora pgvector with HNSW index. The RAG
API Lambda, fronted by API Gateway, embeds user queries with Bedrock,
performs similarity search on pgvector, then generates answers using
Groq/Gemini LLMs. The architecture is detailed below:

\begin{figure}
\centering
\includegraphics{5-Workshop/5.1-Workshop-overview/architecture.png}
\caption{News RAG Pipeline Architecture}
\end{figure}

\hypertarget{aws-services-used}{%
\subsubsection{AWS Services Used}}

\begin{itemize}
\tightlist
\item
  \textbf{Amazon ECS Fargate}: Runs Scrapy SitemapSpider crawler (0.25
  vCPU, 0.5 GB)
\item
  \textbf{Amazon SQS Standard}: Message queue replacing Kafka
  (\textasciitilde\$0/month)
\item
  \textbf{AWS Lambda} (3 functions): Consumer (SQS \(\rightarrow\)
  Aurora), ETL + Bedrock Embed, RAG API
\item
  \textbf{Amazon Aurora Serverless v2}: PostgreSQL 15.4 with pgvector
  extension
\item
  \textbf{Amazon Bedrock}: Titan Embed Text v2 (1024d) for vector
  embeddings
\item
  \textbf{Amazon API Gateway}: REST API frontend for RAG queries
\item
  \textbf{Amazon EventBridge Scheduler}: Daily cron triggers (01:00,
  02:00 UTC)
\item
  \textbf{Amazon ECR}: Docker image registry for Fargate tasks
\item
  \textbf{AWS IAM}: Task execution and Lambda roles with least-privilege
  policies
\item
  \textbf{Amazon CloudWatch}: Centralized logging and monitoring (7-day
  retention)
\end{itemize}

\hypertarget{component-design}{%
\subsubsection{Component Design}}

\begin{itemize}
\tightlist
\item
  \textbf{Crawler (Fargate)}: Scrapy SitemapSpider reads
  sitemap\_news.xml from 3 news sources, parses articles, pushes to SQS.
  Runs daily \textasciitilde30 min.
\item
  \textbf{Queue (SQS)}: Standard queue with DLQ, 14-day retention, 3
  retry limit.
\item
  \textbf{Consumer (Lambda)}: Triggered by SQS, computes SHA256 URL hash
  for dedup, inserts raw articles into Aurora.
\item
  \textbf{ETL + Embed (Lambda)}: Clean HTML, chunk text (500 tokens, 50
  overlap), call Bedrock Titan Embed v2, store 1024d vectors in pgvector
  with HNSW index.
\item
  \textbf{RAG API (Lambda + API Gateway)}: Embed user query via Bedrock,
  cosine similarity search on pgvector, generate answer via Groq/Gemini
  with source citations.
\item
  \textbf{Frontend (Next.js + FastAPI)}: Dashboard with KPI cards,
  charts (Recharts), AI Chat with model selector, Article Explorer,
  Pipeline Monitor.
\end{itemize}

\hypertarget{technical-implementation}{%
\subsubsection{4. Technical
Implementation}}

\hypertarget{implementation-phases}{%
\paragraph{Implementation Phases}}

This project follows 4 phases: - \textbf{Phase 1 - Infrastructure (Weeks
1-2)}: Terraform configuration for VPC, Aurora pgvector, ECS Cluster,
ECR, Lambda, EventBridge, IAM, CloudWatch. Build multi-stage Docker
image for Fargate. - \textbf{Phase 2 - Local Development (Weeks 3-4)}:
Docker Compose environment with PostgreSQL, Qdrant, Kafka. Develop
Scrapy SitemapSpider, Kafka Consumer, ETL pipeline with Star Schema,
SentenceTransformer vectorization. - \textbf{Phase 3 - AWS Production
(Weeks 5-6)}: Deploy Fargate crawler with EventBridge scheduler, Lambda
Consumer with SQS trigger, Lambda ETL with Bedrock Titan Embed v2,
Lambda RAG API with API Gateway, Next.js frontend. - \textbf{Phase 4 -
Testing \& Polish (Weeks 7-8)}: RAGAS evaluation (Faithfulness,
Relevancy, Precision, Recall), CloudWatch dashboards and alerts, Locust
load testing, cost optimization.

\hypertarget{technical-requirements}{%
\paragraph{Technical Requirements}}

\begin{itemize}
\tightlist
\item
  \textbf{Data Pipeline}: Scrapy with SitemapSpider for news crawling,
  Kafka (local) / SQS (AWS) for message queue, PostgreSQL with pgvector
  for storage and vector search.
\item
  \textbf{ETL Pipeline}: HTML cleaning via regex, text chunking at 500
  tokens with 50 overlap, embedding generation via SentenceTransformer
  (local) / Bedrock Titan v2 (AWS).
\item
  \textbf{RAG System}: pgvector HNSW similarity search (cosine
  distance), Groq API (Qwen3-8B, Llama 3.1), Gemini 2.0 Flash fallback,
  structured prompts with source citation.
\item
  \textbf{Infrastructure}: Terraform for all AWS resources, Docker
  multi-stage builds, Lambda deployment packages, EventBridge cron
  scheduling.
\end{itemize}

\hypertarget{timeline-milestones}{%
\subsubsection{5. Timeline \& Milestones}}

\textbf{Project Timeline} - Pre-Internship (Month 0): Planning and AWS
fundamentals study. - Internship (Months 1-3): - Month 1: Infrastructure
setup, local development environment, crawler development. - Month 2:
ETL pipeline, Star Schema, vectorization, RAG API development. - Month
3: AWS production deployment, testing (RAGAS), monitoring, cost
optimization. - Post-Launch: Maintain and extend with additional
features (semantic chunking, hybrid search, topic alerts).

\hypertarget{budget-estimation}{%
\subsubsection{6. Budget Estimation}}

You can use \href{https://calculator.aws/\#/}{AWS Pricing Calculator}
for estimation.

\hypertarget{infrastructure-costs-monthly}{%
\paragraph{Infrastructure Costs
(Monthly)}}

\begin{itemize}
\tightlist
\item
  Aurora Serverless v2 (2 ACU): \textasciitilde\$15-20
\item
  ECS Fargate Crawler (0.25 vCPU, 0.5 GB, 30 min/day):
  \textasciitilde\$0.50
\item
  Lambda (3 functions): \textasciitilde\$2-3
\item
  SQS Standard: \textasciitilde\$0
\item
  API Gateway: \textasciitilde\$0.30
\item
  Bedrock Titan Embed: \textasciitilde\$0.50
\item
  CloudWatch Logs (7-day retention): \textasciitilde\$1-2
\item
  \textbf{Total: \textasciitilde\$21-26/month}
\end{itemize}

\hypertarget{risk-assessment}{%
\subsubsection{7. Risk Assessment}}

\hypertarget{risk-matrix}{%
\paragraph{Risk Matrix}}

\begin{itemize}
\tightlist
\item
  Crawler Blocked by Websites: Medium impact, medium probability
  (mitigate with proper headers, respect robots.txt, DOWNLOAD\_DELAY).
\item
  Bedrock Throttling: Medium impact, low probability (mitigate with
  retry logic, adaptive backoff).
\item
  LLM API Outages (Groq/Gemini): High impact, low probability (mitigate
  with multi-model fallback).
\item
  Cost Overruns: Medium impact, low probability (mitigate with budget
  alerts, serverless right-sizing).
\end{itemize}

\hypertarget{mitigation-strategies}{%
\paragraph{Mitigation Strategies}}

\begin{itemize}
\tightlist
\item
  \textbf{Crawler}: Respect robots.txt, 1s DOWNLOAD\_DELAY, AutoThrottle
  enabled.
\item
  \textbf{Throttling}: Exponential backoff on Bedrock calls,
  max\_attempts=3.
\item
  \textbf{LLM Fallback}: Fallback chain: Groq Qwen3 \(\rightarrow\)
  Llama 3.1 \(\rightarrow\) Gemini 2.0 Flash.
\item
  \textbf{Cost}: AWS Budget alerts at 80\%, Fargate Spot for crawler,
  right-size Lambda memory.
\end{itemize}

\hypertarget{contingency-plans}{%
\paragraph{Contingency Plans}}

\begin{itemize}
\tightlist
\item
  If specific news sites block crawler: Use alternative news sources or
  manual data upload.
\item
  If Bedrock unavailable in region: Use local SentenceTransformer model
  as fallback.
\item
  If costs exceed budget: Reduce Aurora to 1 ACU, increase Lambda
  timeout instead of memory.
\end{itemize}

\hypertarget{expected-outcomes}{%
\subsubsection{8. Expected Outcomes}}

\hypertarget{technical-improvements}{%
\paragraph{Technical Improvements:}}

Automated news aggregation replaces manual searching. AI-powered Q\&A
with source citations saves research time. Scalable serverless
architecture handles increasing news volume.

\hypertarget{long-term-value}{%
\paragraph{Long-term Value}}

\begin{itemize}
\tightlist
\item
  Foundational RAG system for future NLP/AI projects.
\item
  Reusable pipeline components for other domains (tech blogs, research
  papers).
\item
  Hands-on AWS serverless expertise.
\item
  Data foundation (\textasciitilde5000 articles accumulated over
  internship) for future analysis.
\end{itemize}
